\documentclass[11pt]{article}
\usepackage{amsmath, amssymb}
\usepackage{geometry}
\usepackage{hyperref}
\usepackage{graphicx}
\usepackage{booktabs}
\geometry{margin=2.5cm}
\graphicspath{{../plots/}}
\emergencystretch=2em

\newcommand{\Pop}{\hat{P}}
\newcommand{\ket}[1]{|#1\rangle}
\newcommand{\bra}[1]{\langle#1|}
\newcommand{\Oedge}{X_0 Z_1\cdots Z_{L-2} X_{L-1}}
\newcommand{\expv}[1]{\langle #1\rangle}
\newcommand{\Tr}{\operatorname{Tr}}
\newcommand{\cx}{\texttt{cx}}
\newcommand{\chimax}{\chi_{\max}}

\title{Scaling the Kitaev/VQE Simulation to Large $L$\\
\large Exact diagonalization, density matrices, matrix product states, and free fermions}
\author{Seminar notes}
\date{\today}

\begin{document}
\maketitle

\begin{abstract}
Can the Block-4 noise study be pushed from $L=4$ toward $L=16$ (and beyond) on
HPC? This note works through the question honestly. The pipeline hides
\emph{two} exponential objects with very different escape routes. The
\textbf{exact-diagonalization reference} is fully removable: the Kitaev
Hamiltonian is a free-fermion model, so its ground state and observables follow
from a $2L\times 2L$ Bogoliubov--de~Gennes (BdG) matrix in $O(L^3)$ time, exact
to $L\sim$ hundreds. The \textbf{noisy VQE-prepared state} is the hard one:
the \texttt{EfficientSU2} ansatz is a non-Gaussian, essentially universal
circuit, so no classical method escapes exponential cost unconditionally.
Matrix product states (MPS) help \emph{when entanglement is low}, and a direct
probe of the real ansatz shows the bond dimension $\chi$ is bounded
\emph{independent of $L$} at fixed depth ($\chi=8$ for reps${}=3$, all $L$).
But the same probe exposes a fundamental wall: $\chi\approx 2^{\text{reps}}$
exactly, while the fixed-depth ansatz \emph{loses fidelity} with $L$
($0.99\to0.02$ over $L=4\to12$). Keeping the VQE faithful forces
reps${}\sim L$, which sends $\chi$ back to $2^{O(L)}$. MPS is cheap precisely
where the ansatz fails and expensive precisely where it works --- there is no
free lunch, and that \emph{is} the NISQ statement.
\end{abstract}

\section{Two exponential objects}
The current Block-4 evaluation ($\texttt{src/block4.py}$, plots 8--9) contains
two objects that scale exponentially in the number of qubits $L$:
\begin{enumerate}
\item the \textbf{ED reference} --- \texttt{np.linalg.eigh} on the
      $2^L\times 2^L$ Hamiltonian, used for the ideal $\expv{\Oedge}$ curve, the
      parity gap, and the fidelity post-selection;
\item the \textbf{noisy state} --- an exact density matrix $\rho$ of dimension
      $2^L\times 2^L$ (i.e.\ $4^L$ complex entries), produced by
      \texttt{save\_density\_matrix} on the Aer \texttt{density\_matrix} method.
\end{enumerate}
The observables themselves are \emph{not} the problem: \texttt{qubit\_hamiltonian},
\texttt{edge\_string}, and \texttt{parity} are already \texttt{SparsePauliOp}
objects with $O(L)$ Pauli terms. The blow-up comes entirely from calling
\texttt{.to\_matrix()} on them and from \texttt{eigh}. The two objects are
addressed separately below.

\section{The dense wall}
Object 2 sets the practical ceiling. A dense $2^L\times 2^L$ complex\,128 matrix
costs $16\cdot 4^L$ bytes, and the pipeline holds several copies ($\rho$, $H$,
$\Pop$, ED eigenvectors, Aer's internal state):
\begin{center}
\begin{tabular}{rrrl}
\toprule
$L$ & density matrix & $+$ copies & ED \texttt{eigh} cost \\
\midrule
12 & 0.27\,GB & ${\sim}1$\,GB & seconds \\
14 & 4.3\,GB & ${\sim}20$\,GB & minutes \\
15 & 17\,GB & ${\sim}70$\,GB & ${\sim}10$\,min \\
16 & 69\,GB & ${\sim}300$\,GB & hours \\
\bottomrule
\end{tabular}
\end{center}
Worse, \emph{every} VQE cost evaluation is a full $4^L$ density-matrix
simulation, run $\text{maxiter}\times\text{n\_starts}$ times per $L$. Beyond
$L\approx 14$ a single $L$-point becomes hours to days. Even on a large-memory
HPC node the dense path does not reach $L=16$ in feasible wall-clock time; this
is the ``dense wall'' at $L\approx 15$.

\section{Escaping object 1: free fermions}
The Kitaev Hamiltonian is quadratic in fermions. After a Jordan--Wigner map the
qubit form is
\[
H=-\tfrac{\mu}{2}\sum_j Z_j
  +\tfrac{t-\Delta}{2}\sum_j X_jX_{j+1}
  +\tfrac{t+\Delta}{2}\sum_j Y_jY_{j+1},
\]
which is a free-fermion (BdG) model. Its ground state, spectrum, correlation
matrix, and the parity gap $\Delta_P(L)\sim e^{-L/\xi}$ all follow from
diagonalizing a $2L\times 2L$ BdG matrix in $O(L^3)$ time --- \emph{no}
exponential object. The edge operator $\Oedge$ is a product of Majorana
operators, so its ground-state expectation is a Pfaffian of a submatrix of the
covariance matrix (Wick's theorem). This is exactly how Majorana splitting is
computed for $L\sim 1000$ in the literature. \textbf{Conclusion: the ED
reference is removable outright}, exact and cheap to $L\sim$ hundreds. This is
the content of \texttt{src/free\_fermion.py}.

\section{Escaping object 2: trajectories and MPS}
\subsection{Quantum-trajectory statevector sampling}
Noise (a mixed state) does not require the full density matrix. In the
quantum-trajectory picture each shot is a \emph{pure} statevector ($2^L$, not
$4^L$) in which the Kraus operators are sampled stochastically; averaging an
observable over $N_{\text{traj}}$ trajectories reconstructs
$\Tr(O\rho)$ up to a statistical error $\propto 1/\sqrt{N_{\text{traj}}}$.
Memory drops from $4^L$ to $2^L$ ($\sim$1\,MB at $L=16$), at the honest cost of
\emph{error bars} on the noisy curves.

\subsection{Matrix product states}
A general $L$-qubit state is a tensor with $2^L$ amplitudes. An MPS factorizes
it into a chain of small tensors, one per site, linked by bond indices of
dimension $\chi$:
\[
\ket{\psi}=\sum_{\{s\}}A^{s_0}A^{s_1}\cdots A^{s_{L-1}}\ket{s_0\ldots s_{L-1}},
\qquad \text{memory}=O(L\,\chi^2).
\]
Across any cut, $\chi$ equals the number of nonzero Schmidt coefficients --- a
direct measure of \emph{entanglement} across that cut. A gapped 1D ground state
obeys an area law, so its entanglement is bounded independent of $L$; then
$\chi$ is constant and MPS is \emph{linear} in $L$. A critical or volume-law
state needs $\chi\sim 2^{L/2}$ and MPS is exponential again. Aer provides
\texttt{method='matrix\_product\_state'} and supports noise through trajectory
sampling; the empirical test confirms it runs and exposes the per-bond $\chi$.

\subsection{Truncation error, and why the risk is controllable}
A two-qubit gate can double $\chi$ across its bond. To stay cheap one performs
an SVD there and keeps only the largest $\chimax$ Schmidt values. The discarded
weight --- the sum of squared discarded Schmidt coefficients --- \emph{is} the
truncation error: the fraction of the state's norm thrown away. If nothing is
discarded the MPS is \emph{exact}. If $\chimax$ is set below what the state
needs, the simulated state becomes a nearby but different state and
$\expv{\Oedge}$ acquires a systematic bias. Crucially this is a knob with a
readout: Aer reports the discarded weight, and the convergence check ``run at
$\chimax$ and $2\chimax$; if the observable does not move, the result is
effectively exact'' makes the error visible and quantitative rather than
hidden. A truncated number never silently corrupts a result as long as the
discarded weight is logged.

\section{The $\chi$-probe: what the real ansatz actually does}
We measured the \emph{exact} (uncapped) max bond dimension of the real
\texttt{EfficientSU2}(\texttt{ry}, linear, reps) state directly.

\paragraph{$\chi$ is flat in $L$ at fixed depth.}
Worst-of-random-$\theta$ max $\chi$ (reps${}=3$):
\begin{center}
\begin{tabular}{rrrrrr}
\toprule
$L$ & 6 & 10 & 14 & 18 & 22 \\
$\chi$ & 8 & 8 & 8 & 8 & 8 \\
$2^{\lfloor L/2\rfloor}$ (vol.\ law) & 8 & 32 & 128 & 512 & 2048 \\
\bottomrule
\end{tabular}
\end{center}
A reps${}=r$ linear ansatz has a light cone of width ${\sim}r$, so entanglement
across any cut is bounded \emph{independent of $L$}. At reps${}=3$, $\chi$ pins
at $8$ for all $L$, far below the volume-law ceiling, and noise does not inflate
it. \textbf{The truncation risk essentially does not exist for a shallow
ansatz}: cap $\chimax=16$ and the simulation is exact to $L=100$.

\paragraph{$\chi$ grows exponentially in depth.}
At fixed $L=14$:
\begin{center}
\begin{tabular}{rrrrrrrr}
\toprule
reps & 1 & 2 & 3 & 4 & 5 & 6 & 8 \\
$\chi$ & 2 & 4 & 8 & 16 & 32 & 64 & 128 \\
\bottomrule
\end{tabular}
\end{center}
i.e.\ $\chi\approx 2^{\text{reps}}$ (until it saturates the $2^{L/2}$ ceiling).

\paragraph{The fixed-depth ansatz loses the state.}
VQE fidelity to the true ground state at reps${}=3$:
\begin{center}
\begin{tabular}{rrrrrr}
\toprule
$L$ & 4 & 6 & 8 & 10 & 12 \\
fidelity & 0.99 & 0.50 & 0.25 & 0.10 & 0.02 \\
\bottomrule
\end{tabular}
\end{center}
A fixed-depth hardware-efficient ansatz simply cannot represent the topological
ground state as $L$ grows; keeping the fidelity up requires reps${}\sim L$.

\section{The wall is fundamental, not a coding artifact}
Combining the three tables: the regime where MPS is cheap (small reps, constant
$\chi$) is exactly the regime where the VQE \emph{fails} (fidelity $\to 0$); and
the regime where the VQE \emph{works} (reps${}\sim L$) is exactly where
$\chi\sim 2^{\text{reps}}\sim 2^{O(L)}$ blows MPS back up. A state that genuinely
needs a deep circuit to prepare is genuinely expensive to simulate classically
--- MPS buys a smaller exponential base ($2^{\text{reps}}$ vs $4^L$), a few
extra qubits, not unbounded $L$. This is the honest simulation-cost face of
``NISQ does not inherit topological protection.''

\paragraph{The one unconditionally polynomial regime.}
The true Kitaev ground state is Gaussian and area-law: it \emph{has} a small-$\chi$
MPS. The \texttt{ry} ansatz is merely an inefficient, non-Gaussian
parametrization of it. A \textbf{matchgate/Gaussian ansatz} (only
$R_z$-type single-qubit gates $+$ nearest-neighbor matchgate entanglers)
prepares the same state at $O(L)$ depth with bounded $\chi$, and under
\emph{fermion-Gaussian} noise (amplitude damping/loss is Gaussian for fermions;
depolarizing is not) the entire simulation reduces to covariance-matrix algebra,
$O(L^3)$, exact, $L\to$ thousands. The catch: this answers a cleaner but
\emph{different} question --- it is no longer the generic hardware-efficient
NISQ story.

\section{Recommended architecture}
\begin{center}
\begin{tabular}{lll}
\toprule
piece & scalable replacement & exact? \\
\midrule
ideal $\expv{\Oedge}$, gap & free-fermion BdG ($L^3$) & exact, $L\to$ hundreds \\
noisy VQE $\expv{\Oedge}$ & MPS $+$ noise trajectories ($L\chi^2$) & exact if $\chi$ bounded \\
unbounded $L$ & matchgate ansatz $+$ Gaussian noise ($L^3$) & exact, different model \\
\bottomrule
\end{tabular}
\end{center}
The concrete plan: build \texttt{src/free\_fermion.py} (an unconditional win)
and pair it with a \emph{fixed-reps} MPS sweep that makes the ansatz's
expressibility collapse visible against the exact ideal curve. That is honest,
cheap, reaches large $L$, and turns the wall itself into the scientific point.
Validation anchor: at $L\le 12$, where the exact density-matrix path still runs,
cross-check the MPS-trajectory and free-fermion values against the exact
numbers before extrapolating.

\paragraph{Backends for the noisy curve.} A device-calibrated sweep needs a
snapshot whose coupling map embeds a length-$L$ line so the linear-entanglement
\cx{} pairs sit on real edges with no SWAP inflation. Named 5-qubit snapshots
(\texttt{fake\_manila}, etc.) cap at $L=5$; \texttt{FakeGuadalupeV2} is 16 qubits
but its heavy-hex graph only contains a 13-node line. The clean choices are the
27-qubit heavy-hex snapshots (\texttt{FakeCairoV2}/\texttt{FakeKolkataV2}/%
\texttt{FakeHanoiV2}, 21-node line) or a 20-qubit device (18-node line).

\paragraph{Reproduce.} $\chi$-probe:
\texttt{python scratch/chi\_probe.py} (measures $\chi$ vs $L$ and vs reps against
the volume-law ceiling). Free-fermion reference: \texttt{src/free\_fermion.py}
(BdG ground state, edge string, parity gap; validated against \texttt{eigh} at
$L\le 12$).

\end{document}
